%----------------------------------------------------------------------------------------
%	PACKAGES AND THEMES
%----------------------------------------------------------------------------------------
\documentclass[aspectratio=169,xcolor=dvipsnames]{beamer}
\usetheme{SimplePlus}

\usepackage{hyperref}
\usepackage{graphicx} % Allows including images
\usepackage{booktabs} % Allows the use of \toprule, \midrule and \bottomrule in tables

%----------------------------------------------------------------------------------------
%	TITLE PAGE
%----------------------------------------------------------------------------------------

\title[short title]{The End of the Crypto-Diversification Myth} % The short title appears at the bottom of every slide, the full title is only on the title page
\subtitle{Advanced Financial Theory - MsQEF}

\author[Pin-Yen] {Antiga Ibadova, Reha Tuncer}

\institute[Uni Lu] % Your institution as it will appear on the bottom of every slide, may be shorthand to save space
{
    Faculty of Law, Economics and Finance (FDEF) \\
    University of Luxembourg % Your institution for the title page
}
\date{12 December 2024} % Date, can be changed to a custom date


%----------------------------------------------------------------------------------------
%	PRESENTATION SLIDES
%----------------------------------------------------------------------------------------

\begin{document}

\begin{frame}
    % Print the title page as the first slide
    \titlepage
\end{frame}

\begin{frame}{The Crypto-Diversification Puzzle}
    \begin{figure}
    \centering
    \includegraphics[width=0.7\linewidth]{fig1.png}
    \end{figure}
     \begin{itemize}
    \item Bitcoin-S\&P500 correlation jumps from 0 to 0.6 in March 2020 and remains positive
    \item Challenging the stock market diversification with Bitcoin narrative
    \item Institutional investors out of the picture (at the time) with uncorrelated fundamentals  
    \end{itemize}
\end{frame}

\begin{frame}{Research Question}
    \begin{itemize}
    \item  What drives the Bitcoin and stock cross-asset price correlation since March 2020?
\item \textbf{Mechanism:} Retail investors' trading habits
\item Theory driven model of retail investors with testable hypotheses
\item Empirical evidence from unique dataset: A Swiss bank's crypto trading clients
    \end{itemize}
\end{frame}

\begin{frame}{Theoretical Framework}
   \begin{itemize}
\item Extension of the Kyle (1985) model to two assets, crypto and stock
    \begin{enumerate}
\item Informed ``insider'' and uninformed ``noise'' traders with market makers who cannot distinguish between these two
\item Price correlation between markets come from spillovers between the two markets
\end{enumerate}

\item \textbf{Assumptions}
    \begin{enumerate}
\item Uncorrelated fundamentals: No mechanism links Bitcoin and stock market fluctuations
\item Segmented market-making: Sum of trades in one market cannot be observed by the other (reasonable for 2020)
\item Correlated uninformed investor trading: Sentiment ($+$) or wealth effects ($-$) 
\end{enumerate}
\item \textbf{Prediction \textcolor{red}{work here}} 
$Cov(p_1,p_2) > 0 \iff \rho > 0$ (retail correlation drives price correlation)
\end{itemize}
\end{frame}

\begin{frame}{Data}
    \begin{itemize}
\item Opening with the correlation graph (the one provided)
\item Key question: Why did Bitcoin-S\&P500 correlation jump from 0 to 0.6 in March 2020?
\item Challenge to the "crypto as diversification" narrative
\item Motivation: Policy implications of crypto-equity market linkages
    \end{itemize}
\end{frame}

\begin{frame}{Identification Strategy}
    \begin{itemize}
\item Opening with the correlation graph (the one provided)
\item Key question: Why did Bitcoin-S\&P500 correlation jump from 0 to 0.6 in March 2020?
\item Challenge to the "crypto as diversification" narrative
\item Motivation: Policy implications of crypto-equity market linkages
    \end{itemize}
\end{frame}

\begin{frame}{Data}
    \begin{itemize}
\item Opening with the correlation graph (the one provided)
\item Key question: Why did Bitcoin-S\&P500 correlation jump from 0 to 0.6 in March 2020?
\item Challenge to the "crypto as diversification" narrative
\item Motivation: Policy implications of crypto-equity market linkages
    \end{itemize}
\end{frame}

\begin{frame}{Results 1}
    \begin{itemize}
\item Opening with the correlation graph (the one provided)
\item Key question: Why did Bitcoin-S\&P500 correlation jump from 0 to 0.6 in March 2020?
\item Challenge to the "crypto as diversification" narrative
\item Motivation: Policy implications of crypto-equity market linkages
    \end{itemize}
\end{frame}

\begin{frame}{Results 2}
    \begin{itemize}
\item Opening with the correlation graph (the one provided)
\item Key question: Why did Bitcoin-S\&P500 correlation jump from 0 to 0.6 in March 2020?
\item Challenge to the "crypto as diversification" narrative
\item Motivation: Policy implications of crypto-equity market linkages
    \end{itemize}
\end{frame}

\begin{frame}{Results 3}
    \begin{itemize}
\item Opening with the correlation graph (the one provided)
\item Key question: Why did Bitcoin-S\&P500 correlation jump from 0 to 0.6 in March 2020?
\item Challenge to the "crypto as diversification" narrative
\item Motivation: Policy implications of crypto-equity market linkages
    \end{itemize}
\end{frame}

\begin{frame}{Contribution}
    % Throughout your presentation, if you choose to use \section{} and \subsection{} commands, these will automatically be printed on this slide as an overview of your presentation
    \tableofcontents
\end{frame}

\begin{frame}{Implications and Current Relevance}
    % Throughout your presentation, if you choose to use \section{} and \subsection{} commands, these will automatically be printed on this slide as an overview of your presentation
    \tableofcontents
\end{frame}

%------------------------------------------------
\section{First Section}
%------------------------------------------------

\begin{frame}{Bullet Points}
    \begin{itemize}
        \item Lorem ipsum dolor sit amet, consectetur adipiscing elit
        \item Aliquam blandit faucibus nisi, sit amet dapibus enim tempus eu
        \item Nulla commodo, erat quis gravida posuere, elit lacus lobortis est, quis porttitor odio mauris at libero
        \item Nam cursus est eget velit posuere pellentesque
        \item Vestibulum faucibus velit a augue condimentum quis convallis nulla gravida
    \end{itemize}
\end{frame}

%------------------------------------------------

\begin{frame}{Blocks of Highlighted Text}
    In this slide, some important text will be \alert{highlighted} because it's important. Please, don't abuse it.

    \begin{block}{Block}
        Sample text
    \end{block}

    \begin{alertblock}{Alertblock}
        Sample text in red box
    \end{alertblock}

    \begin{examples}
        Sample text in green box. The title of the block is ``Examples".
    \end{examples}
\end{frame}

%------------------------------------------------

\begin{frame}{Multiple Columns}
    \begin{columns}[c] % The "c" option specifies centered vertical alignment while the "t" option is used for top vertical alignment

        \column{.45\textwidth} % Left column and width
        \textbf{Heading}
        \begin{enumerate}
            \item Statement
            \item Explanation
            \item Example
        \end{enumerate}

        \column{.5\textwidth} % Right column and width
        Lorem ipsum dolor sit amet, consectetur adipiscing elit. Integer lectus nisl, ultricies in feugiat rutrum, porttitor sit amet augue. Aliquam ut tortor mauris. Sed volutpat ante purus, quis accumsan dolor.

    \end{columns}
\end{frame}

%------------------------------------------------
\section{Second Section}
%------------------------------------------------

\begin{frame}{Table}
    \begin{table}
        \begin{tabular}{l l l}
            \toprule
            \textbf{Treatments} & \textbf{Response 1} & \textbf{Response 2} \\
            \midrule
            Treatment 1         & 0.0003262           & 0.562               \\
            Treatment 2         & 0.0015681           & 0.910               \\
            Treatment 3         & 0.0009271           & 0.296               \\
            \bottomrule
        \end{tabular}
        \caption{Table caption}
    \end{table}
\end{frame}

%------------------------------------------------

\begin{frame}{Theorem}
    \begin{theorem}[Mass--energy equivalence]
        $E = mc^2$
    \end{theorem}
\end{frame}

%------------------------------------------------

\begin{frame}{Figure}
    Uncomment the code on this slide to include your own image from the same directory as the template .TeX file.
    %\begin{figure}
    %\includegraphics[width=0.8\linewidth]{test}
    %\end{figure}
\end{frame}

%------------------------------------------------

\begin{frame}[fragile] % Need to use the fragile option when verbatim is used in the slide
    \frametitle{Citation}
    An example of the \verb|\cite| command to cite within the presentation:\\~

    This statement requires citation \cite{p1}.
\end{frame}

%------------------------------------------------

\begin{frame}{References}
    % Beamer does not support BibTeX so references must be inserted manually as below
    \footnotesize{
        \begin{thebibliography}{99}
            \bibitem[Smith, 2012]{p1} John Smith (2012)
            \newblock Title of the publication
            \newblock \emph{Journal Name} 12(3), 45 -- 678.
        \end{thebibliography}
    }
\end{frame}

%------------------------------------------------

\begin{frame}
    \Huge{\centerline{\textbf{The End}}}
\end{frame}

%----------------------------------------------------------------------------------------

\end{document}